\subsection{5月}

\subsubsection{5月15日 周五}

\begin{enumerate}
  \item \textbf{SPI 模块挂载}：在 AXI-lite 总线上完成 SPI 控制器的挂载，并编写了对应的 C 驱动程序。
  \item \textbf{UART Monitor}：完成了一个仅用于仿真环境的 UART monitor 模块。
  \item \textbf{Hermes 个人系统构建}：搭建了基于 Hermes Agent 的个人工作流系统，配置了若干 cron 定时任务，分别用于同步 Nix 知识库以及 Hermes 记忆与配置。
\end{enumerate}

昨日被导师陈鑫严厉批评，指摘我不把 \emph{deadline} 当回事，甚至放出狠话「是不是找打」。警醒之至，当引以为戒。时间管理不可再松弛，任务节点必须死守。


\subsubsection{5月16日 周六}

今天花了一天研究三个 ADC 校准算法的文档资料，借助 AI 辅助理解算法原理和实现思路。后续计划写个 C 程序验证校准逻辑，再制作一份 PPT 用于汇报展示——毕竟没有实际硬件环境，仿真验证就够了。周一还得把代码导入到内网开发环境。

\begin{itemize}
  \item \textbf{ADC 校准算法研究}：仔细阅读了三份校准文档，理解各算法的数学原理和适用场景。AI 辅助翻译和解释公式推导，帮助厘清算法间的差异。
  \item \textbf{后续计划}：编写 C 验证程序 + 制作 PPT，周一导入内网。
  \item 李一航计划明天去看房子，准备租房事宜。
  \item 月底需要审车（车辆年检），驾照 8 月份也到期要换证了——记得提前预约体检和照相。
  \item 天气盛美 ☀️🌸，就是风有点大。
\end{itemize}

今天感受：\\[3pt]
早上老陈 9 点多就开始催进度 😅，紧迫感拉满。不过研究下来发现 ADC 校准的思路逐渐清晰，心里还算有底。得抓紧时间出代码和 PPT，千万不能再拖了。

